% =====================================================================
%  05-criptografia.tex - Resultados criptograficos
% =====================================================================

\section{Resultados criptográficos}\label{sec:pruebas-cripto}

\subsection{Vectores oficiales}\label{sec:pruebas-vectores}

Todos los vectores oficiales adoptados están en verde. La \cref{tab:pruebas-vectores} los resume. Para SHA-256 se reproducen los cuatro vectores de ejemplo de FIPS 180-4 \cite{nist2015fips1804}: la cadena vacía, el mensaje «abc», el mensaje de dos bloques y un millón de caracteres «a» ---este último consumido por flujo, lo que ejercita a la vez la interfaz de streaming del adaptador---. El vector «abc» se reutiliza además de extremo a extremo: la prueba del subcomando \texttt{crypto hash} verifica que el binario imprime exactamente \texttt{ba7816bf}\ldots{} ante ese contenido. Para AES-256-GCM se adoptaron los casos de prueba TC13 a TC16 de la especificación del modo GCM de McGrew y Viega \cite{mcgrewviega2004}, la base de NIST SP 800-38D \cite{dworkin2007}; los mismos valores figuran en las suites del programa de validación CAVP del NIST. El adaptador expone una operación de sellado con nonce inyectable exclusivamente para reproducir los vectores; las cuatro pruebas de cifrado exigen el criptograma y el tag exactos, dos de descifrado exigen la autenticación y el texto claro exactos, y una prueba de respuesta conocida negativa voltea el último byte del tag y exige el rechazo autenticado.

\begin{table}[H]
  \centering
  \caption{Vectores oficiales reproducidos por la suite de pruebas.}
  \label{tab:pruebas-vectores}
  \begin{tabularx}{\linewidth}{@{}L{4.4cm} X C{1.9cm}@{}}
    \toprule
    \textbf{Estándar y fuente} & \textbf{Vectores reproducidos} & \textbf{Pruebas} \\
    \midrule
    FIPS 180-4, SHA-256 \cite{nist2015fips1804} & Cadena vacía (\texttt{e3b0c442}\ldots), «abc» (\texttt{ba7816bf}\ldots), mensaje de dos bloques (\texttt{248d6a61}\ldots) y un millón de «a» por flujo (\texttt{cdc76e5c}\ldots). & 4 \\
    \addlinespace
    Casos TC13--TC16 del modo GCM \cite{mcgrewviega2004,dworkin2007} & AES-256-GCM: cifrado de los cuatro casos con criptograma y tag exactos; descifrado autenticado de TC14 y TC16; caso negativo con el último byte del tag volteado. & 7 \\
    \addlinespace
    RFC 6238, apéndice B \cite{mraihi2011} & Los seis vectores TOTP SHA-1 de ocho dígitos ($T$ desde 59 hasta 20\,000\,000\,000, secreto ASCII de referencia). & 1 (6 vectores) \\
    \addlinespace
    RFC 4226, apéndice D \cite{mraihi2005} & Los diez vectores HOTP de seis dígitos (contadores 0 a 9), como verificación cruzada del conteo de dígitos de producción. & 1 (10 vectores) \\
    \addlinespace
    CRC-32 del formato ZIP & Valor de comprobación estándar: \texttt{crc32} de \texttt{"123456789"} igual a \texttt{0xCBF43926}, y \texttt{crc32} de la cadena vacía igual a 0. & 1 \\
    \bottomrule
  \end{tabularx}
\end{table}

Para la autenticación multifactor se reproducen íntegros los dos apéndices normativos: los seis vectores TOTP de RFC 6238 con SHA-1 y ocho dígitos, y los diez vectores HOTP de RFC 4226, que cruzan el conteo de seis dígitos usado en producción. Para la firma digital RSA-3072 con padding PKCS\#1 v1.5 no se adoptaron vectores oficiales ---los publicados fijan claves de otros tamaños---; su lugar lo ocupan dos verificaciones equivalentes: el determinismo del esquema, fijado por la prueba de que las codificaciones PKCS\#1 y PKCS\#8 de una misma clave producen firmas byte a byte idénticas, y la interoperabilidad bidireccional con OpenSSL que documenta la \cref{sec:pruebas-interop}.

\subsection{Frescura de nonces}\label{sec:pruebas-nonces}

El criterio del objetivo OE-4 exige cero reutilizaciones de nonce en una suite de 10\,000 cifrados. La prueba \texttt{nonces\_are\_distinct\_across\_ten\_thousand\_seals} ejecuta 10\,000 operaciones de sellado consecutivas sobre el cifrador real y registra el nonce de 96 bits de cada una en un conjunto: la prueba falla si algún nonce se repite. El resultado es de 10\,000 nonces distintos en 10\,000 operaciones, consistente con la generación por CSPRNG fresco en cada sellado que estableció el diseño (\cref{sec:cifrado-reposo}). La misma familia de garantías de frescura está fijada por pruebas en el resto del sistema: las dos DEK generadas en la prueba son distintas, dos enrolamientos TOTP extraen secretos distintos, cada hash de contraseña usa una sal aleatoria nueva ---dos hashes de la misma contraseña difieren y ambos verifican---, los ocho códigos de recuperación de un lote no se repiten, y ocho solicitudes concurrentes a la autoridad de sellado local producen ocho números de serie distintos.

\subsection{Detección de mutaciones}\label{sec:pruebas-mutacion}

Las pruebas de mutación alteran un artefacto criptográfico real ---un byte o un bit a la vez--- y exigen el rechazo. El resultado global es la detección del 100\,\% de las mutaciones ensayadas, en todos los servicios donde se aplicaron. La pieza central es la prueba exhaustiva del cifrado autenticado, \texttt{any\_single\_byte\_flip\_breaks\_authentication}: recorre en bucle \emph{cada} byte del payload sellado ---nonce, criptograma y tag, uno por uno--- y verifica que voltear cualquiera de ellos hace fallar la apertura con el error opaco de autenticación; para este servicio la afirmación de detección total no es muestral sino exhaustiva sobre todas las posiciones del artefacto. Alrededor de ese núcleo, la suite ensaya mutaciones sobre cada tipo de evidencia:

\begin{itemize}
  \item \textbf{Firma digital:} un byte alterado del documento firmado, o un byte alterado de la firma, produce el rechazo tipado de no correspondencia; el mismo caso se ensaya de extremo a extremo, donde el componente de firma de la verificación integral falla con su causa y el veredicto pasa a no válido.
  \item \textbf{Certificado:} corromper el último byte DER de un certificado ---su firma--- hace que el validador reporte emisor no confiable aunque la estructura siga parseando; la prueba fija que la confianza depende de la verificación criptográfica y no del parseo.
  \item \textbf{Token de sellado:} un token DER truncado o sustituido por bytes aleatorios se diagnostica como malformado; un token legítimo evaluado contra un ancla ajena queda como no confiable, conservando el diagnóstico de la herramienta.
  \item \textbf{Cadena de auditoría:} \texttt{any\_single\_corruption\_is\_detected\_at\_its\_index} corrompe un bit de un byte de cualquier campo ---valor de cadena, secuencia, actor, acción o recurso--- de cualquier entrada de bitácoras generadas aleatoriamente, y exige que la verificación se rompa exactamente en el índice alterado; el mismo comportamiento se fija a mano sobre el archivo JSONL y de extremo a extremo por el subcomando \texttt{audit verify-chain}.
  \item \textbf{Paquete cifrado y paquete de evidencia:} un byte alterado del archivo \texttt{DVLT1} produce el rechazo autenticado en el subcomando \texttt{vault decrypt}; un byte alterado de un miembro del ZIP lo detecta la herramienta independiente \texttt{unzip} por el CRC; un documento corrompido dentro del paquete extraído es rechazado por los comandos \texttt{openssl} del propio instructivo; y la exportación se rehúsa a escribir un paquete cuya firma fue alterada, nombrando el componente fallido.
\end{itemize}

Este conjunto aporta evidencia del criterio de detección del objetivo OE-4 para las mutaciones y artefactos ensayados y materializa la mitigación comprometida en el análisis de riesgos para el compromiso del esquema de cifrado (RTC-02, \cref{tab:rtc-02}).

\subsection{Propiedades verificadas por generación aleatoria de casos}\label{sec:pruebas-propiedades}

Cinco invariantes se verifican con \texttt{proptest} sobre cientos de entradas generadas aleatoriamente por corrida, en la tradición de las herramientas de property-based testing \cite{claessen2000}:

\begin{enumerate}
  \item \textbf{Round-trip del cifrado autenticado:} para toda llave de 32 bytes, todo AAD de hasta 63 bytes y todo texto claro de hasta 1\,023 bytes, abrir lo sellado devuelve el texto claro original.
  \item \textbf{Determinismo de SHA-256:} para entradas arbitrarias de hasta 4\,095 bytes, el digest es idéntico entre corridas y entre la interfaz de bloque y la de flujo.
  \item \textbf{Sensibilidad de SHA-256:} voltear cualquier bit de cualquier byte de la entrada cambia el digest.
  \item \textbf{Validez de la cadena de auditoría:} cualquier secuencia aleatoria de 0 a 11 eventos anexados produce una cadena que verifica como válida.
  \item \textbf{Detección puntual de corrupción:} una corrupción de un bit en cualquier campo de cualquier entrada de la bitácora se detecta exactamente en su índice.
\end{enumerate}

Las dos últimas propiedades son el complemento probabilístico de las pruebas de mutación de la \cref{sec:pruebas-mutacion}: no fijan un caso concreto sino el invariante completo del que esos casos son instancias.

\subsection{Concurrencia}\label{sec:pruebas-concurrencia}

Los dos adaptadores que comparten recursos del sistema de archivos entre procesos ---la bitácora encadenada y la autoridad de sellado local--- se probaron bajo concurrencia real, no simulada. Ocho hilos que anexan cuatro eventos cada uno sobre el mismo archivo de bitácora, liberados simultáneamente con una barrera, producen 32 entradas con secuencias únicas y una sola cadena que verifica como válida: el lock exclusivo entre procesos serializa los escritores en lugar de permitir que bifurquen la cadena, y la misma garantía se fija para instancias del adaptador en procesos separados que escriben por turnos. En la autoridad de sellado, ocho solicitudes simultáneas producen ocho tokens con números de serie todos distintos ---la condición de carrera sobre el archivo de serie de OpenSSL que el adaptador debe serializar (\cref{sec:impl-s5-sellado})---, cuatro verificaciones simultáneas del mismo token concluyen todas válidas gracias al directorio de trabajo privado de cada una, y ninguna operación deja archivos de trabajo residuales en el directorio de la autoridad.

\subsection{Verificaci\'on del portal administrativo CA}
\label{sec:pruebas-portal-ca}

La aceptaci\'on del servicio descrito en \cref{sec:impl-portal-ca} separ\'o
enrolamiento, recorrido HTTP, navegador y recuperaci\'on. Seis casos del alta
administrativa comprobaron entrega privada y confirmaci\'on TOTP antes de la
inserci\'on at\'omica. Incluyeron factor incorrecto, archivo de entrega modificado,
contrase\~nas distintas, fin de entrada, rutas existentes o enlazadas y cancelaci\'on
con restauraci\'on del eco de terminal. La repetici\'on de un alta ya confirmada
no produjo otra cuenta ni otra entrega. La evidencia est\'a expresada en
\rutarepo{crates/bin/tests/ca_enrollment.rs}; no acredita un alta del operador real
en el servidor de destino por s\'i sola.

El recorrido HTTP nativo aprob\'o con PostgreSQL, Redis, el binario CA, el worker
y OpenSSL reales en un entorno desechable. Despu\'es del login y MFA, inspeccion\'o
un CSR de clave externa, emiti\'o ambos perfiles, repiti\'o la operaci\'on exacta,
consult\'o inventario e historial y descarg\'o el certificado identificado por
serie y huella. La revocaci\'on repetida conserv\'o un solo recibo y qued\'o
pendiente de publicaci\'on, sin adelantar por s\'i misma la CRL. Peticiones sin
sesi\'on, un desaf\'io usado como token, autoridad del despacho y campos de actor
o rutas introducidos por el cliente se rechazaron sin nueva mutaci\'on PKI.
Este ensayo no ejecut\'o una parada de los servicios de producci\'on.

Dos recorridos de navegador, a 1440 y 390 p\'ixeles, aprobaron con transporte
HTTP controlado. Ejercitaron revisi\'on y confirmaci\'on de la solicitud,
respuesta incierta, reanudaci\'on desde el historial, revocaci\'on local y
publicaci\'on. Ese control permite comprobar la conservaci\'on del UUID y de
la intenci\'on original ante respuestas tard\'ias; no sustituye la criptograf\'ia
y persistencia del recorrido HTTP nativo. Los componentes visuales compartidos
no implican que la sesi\'on CA y la del despacho sean intercambiables.

Un grupo focal de 55 casos comprob\'o guardas, mantenimiento CRL y respaldo.
La incorporaci\'on de la guarda real al runtime controlado reprodujo que el
journal de publicaci\'on pendiente imped\'ia el arranque despu\'es del commit.
La correcci\'on mantuvo las barreras CRL y de restauraci\'on en el arranque y
reserv\'o el rechazo de otras operaciones pendientes a las mutaciones. La
captura gen\'erica y la restauraci\'on continuaron bloqueadas ante esas operaciones.
Las pruebas del tar exigieron los tres archivos de recuperaci\'on de la
publicaci\'on, rechazaron ausencias y enlaces y conservaron el alcance previo
de una captura sin journal del portal. Los dobles de servicios de este grupo no
demuestran por s\'i mismos cierre ni reapertura de unidades desplegadas.

El ensayo nativo de restauraci\'on SQL aprob\'o en 12.37~s. Volc\'o el esquema
CA poblado, retir\'o ese esquema y lo restaur\'o, contrastando los datos de cuenta
y auditor\'ia con la captura previa. Un c\'odigo de recuperaci\'on consumido antes
del respaldo permaneci\'o consumido. La sesi\'on previamente cerrada continu\'o
rechazada y una nueva autenticaci\'on multifactor, con otro c\'odigo disponible,
recuper\'o el mismo UUID y correo. La cadena de auditor\'ia se verific\'o y la ruta
de expedientes sigui\'o ausente del servicio CA, con el estado del despacho
conservado. El caso reside en \rutarepo{crates/bin/tests/ca_restore_service.rs};
no restaur\'o simult\'aneamente el RDB ni el material PKI.

Un ensayo separado con Redis nativo aprob\'o en 5.924~s. Restaur\'o el RDB con
sesiones y desaf\'ios CA, invalid\'o esas capacidades junto a las del despacho y
conserv\'o tipos, valores y expiraciones absolutas de cuotas y marcas TOTP.
Despu\'es de persistir el estado depurado en AOF, reinici\'o con el RDB anterior
presente y comprob\'o que las capacidades eliminadas no reaparecieran. La prueba
\rutarepo{scripts/tests/test_deployment_certificate_login_snapshot.py} emple\'o
procesos Redis propios; PostgreSQL y la observaci\'on de systemd permanecieron
controlados. Las dos restauraciones conservan, por tanto, fronteras distintas.

Un caso adicional extrajo el tar real de una publicaci\'on preparada y
conserv\'o byte por byte PKI, journals y configuraci\'on con una KEK CA distinta.
La reanudaci\'on del mismo UUID us\'o el candidato conservado sin generar otra
CRL ni cambiar la serie. SQL y servicios fueron controlados en ese caso; no
convierte las dos restauraciones nativas separadas en un ensayo de host completo.

Una correlaci\'on nativa posterior aprob\'o en 37.955~s el uso de certificados
emitidos por la CA en ambos consumidores: acceso Owner con RSA y MFA, y firma
documental con sello y paquete verificable. Una CRL sucesora rechaz\'o usos
nuevos y conserv\'o recibos, auditor\'ia y paquete hist\'orico. El ensayo emple\'o
identidades sint\'eticas y supervisi\'on controlada; no acredita mantenimiento
del host de producci\'on.

El despacho se activ\'o en la release \texttt{v0.1.8} el 13 de octubre de 2026.
El 14 de octubre se completaron el enrolamiento real del operador CA, su
configuraci\'on privada y el arranque de ambas unidades independientes. El proxy
HTTPS propio se abri\'o sin modificar el del despacho. El navegador comprob\'o
contrase\~na y TOTP y mostr\'o la identidad configurada. El acceso Owner por
certificado permaneci\'o deshabilitado, conservando sus credenciales anteriores.

La prueba p\'ublica de emisi\'on encontr\'o una solicitud con nombre com\'un,
pero sin pa\'is ni organizaci\'on. La revisi\'on previa la admiti\'o y OpenSSL
la rechaz\'o despu\'es de guardar la intenci\'on. No cambiaron el \'indice ni
el contador y no se gener\'o un certificado; reanudar repiti\'o el fallo. La
operaci\'on se conserv\'o pendiente y motiv\'o la admisi\'on anticipada del sujeto
y el rechazo terminal limitado al estado intacto descritos en
\cref{sec:impl-portal-ca}.

La reparaci\'on local aprob\'o cuatro pruebas Rust, 60 pruebas Python con
inspector nativo, 23 pruebas Node y un recorrido de navegador. Se verificaron
atributos ausentes, diferentes, repetidos y con distinto tipo ASN.1; un fallo
gen\'erico o material alterado mantuvo la operaci\'on pendiente. El recibo de
rechazo conserv\'o la intenci\'on original, y la interfaz no ofreci\'o descargar
un certificado inexistente. Esta reparaci\'on a\'un no est\'a desplegada ni ha
conciliado el incidente p\'ublico. Permanecen pendientes las mutaciones del
portal instalado, el recorrido real Owner y la recuperaci\'on integral del host.
El inventario de evidencia y sus l\'imites se conserva en
\rutarepo{docs/ca/verification.md}.

